% concatenated from BPGL1splittings.tex, BPGL1splittings/preamble.tex
\documentclass{amsart}
\usepackage[margin=1.5in]{geometry}
%%%%% This document is a shared preamble file 

\setcounter{tocdepth}{1}

\parskip 0.7pc
\parindent 0pt

%\usepackage[margin=1in]{geometry}

\allowdisplaybreaks[1]

\usepackage{url} % Native XeTeX URL handling; original UrlBreaks retained
\def\UrlBreaks{\do\/\do-}

\usepackage{color}
\usepackage{tikz-cd}
\usetikzlibrary{arrows}
\usepackage{amsmath,amssymb}
\usepackage{graphicx}
\usepackage{mathrsfs}
\usepackage{mathabx}
\usepackage[all]{xy}
\usepackage{adjustbox}
\usepackage[normalem]{ulem}
%\usepackage{sseq}
% Unused spectral-sequence drawing package; all original charts are PNG assets.
\usepackage{tcolorbox}
\usepackage{xfrac}
\usepackage{soul}
\usepackage{varwidth}

\usepackage{hyperref}
\usepackage[capitalise, noabbrev]{cleveref}
\usepackage[T1]{fontenc}
\usepackage[utf8]{inputenc}
\usepackage{caption}
\usepackage{subcaption}
\usepackage{multicol}

%%%%%%% for numbered theorems %%%%%%%%%
\newtheorem{theorem}{Theorem}[section]
\newtheorem{cor}[theorem]{Corollary}
\newtheorem{lemma}[theorem]{Lemma}
\newtheorem{proposition}[theorem]{Proposition}
\newtheorem{obs}[theorem]{Observation}
\newtheorem{remark}[theorem]{Remark}
\newtheorem{definition}[theorem]{Definition}
\newtheorem{example}[theorem]{Example}
\newtheorem{observation}[theorem]{Observation}

\newtheorem*{thm*}{Theorem}
\newtheorem*{cor*}{Corollary}
\newtheorem*{lem*}{Lemma}
\newtheorem*{prop*}{Proposition}

\theoremstyle{plain}
\newtheorem{rmk}[theorem]{Remark}

%%%%%%%%%%%%%%%%%%%%%%%%%%%%%%%%%%%%%%%
\numberwithin{theorem}{section}
\numberwithin{equation}{section}
\numberwithin{figure}{section}
\DeclareMathOperator{\Ext}{Ext}
\DeclareMathOperator{\Tor}{Tor}
\DeclareMathOperator{\Hom}{Hom}
\DeclareMathOperator{\Prim}{Prim}
\DeclareMathOperator{\aut}{Aut}
\DeclareMathOperator{\im}{im}
\DeclareMathOperator{\Sta}{Sta}
\DeclareMathOperator{\Map}{Map}
\DeclareMathOperator*{\holim}{holim}
\DeclareMathOperator*{\hocolim}{hocolim}
\DeclareMathOperator*{\colim}{colim}
\DeclareMathOperator*{\Tot}{Tot}
\DeclareMathOperator{\Spf}{Spf}
\DeclareMathOperator{\Aut}{Aut}
\DeclareMathOperator{\Spec}{Spec}
\DeclareMathOperator{\Proj}{Proj}
\DeclareMathOperator{\wt}{wt}
\DeclareMathOperator{\Sq}{Sq}
\DeclareMathOperator{\gr}{gr}

\DeclareMathOperator{\sq}{Sq}
\newcommand{\xib}{\zeta}


\newcommand{\E}[2]{\prescript{#1}{#2}{E}}

\newcommand{\mB}{{\mathbb B}}
\newcommand{\mC}{{\mathbb C}}
\newcommand{\mD}{{\mathbb D}}
\newcommand{\mE}{{\mathbb E}}
\newcommand{\mF}{{\mathbb F}}
\newcommand{\mG}{{\mathbb G}}
\newcommand{\mK}{{\mathbb K}}
\newcommand{\mM}{{\mathbb M}}
\newcommand{\mN}{{\mathbb N}}
\newcommand{\mQ}{{\mathbb Q}}
\newcommand{\mR}{{\mathbb R}}
\newcommand{\mS}{{\mathbb S}}
\newcommand{\mT}{{\mathbb T}}
\newcommand{\mZ}{{\mathbb Z}}

\newcommand{\cA}{{\mathcal A}}
\newcommand{\cB}{{\mathcal B}}
\newcommand{\cC}{{\mathcal C}}
\newcommand{\cE}{{\mathcal E}}
\newcommand{\cF}{{\mathcal F}}
\newcommand{\cG}{{\mathcal G}}
\newcommand{\cI}{{\mathcal I}}
\newcommand{\cJ}{{\mathcal J}}
\newcommand{\cK}{{\mathcal K}}
\newcommand{\cL}{{\mathcal L}}
\newcommand{\cM}{{\mathcal M}}
\newcommand{\cN}{{\mathcal N}}
\newcommand{\cO}{{\mathcal O}}
\newcommand{\cP}{{\mathcal P}}
\newcommand{\cQ}{{\mathcal Q}}
\newcommand{\cS}{{\mathcal S}}
\newcommand{\cT}{{\mathcal T}}
\newcommand{\cU}{{\mathcal U}}
\newcommand{\cV}{{\mathcal V}}
\newcommand{\cW}{{\mathcal W}}
\newcommand{\cX}{{\mathcal X}}
\newcommand{\cZ}{{\mathcal Z}}

\newcommand{\umF}{\underline{\mF}}
\newcommand{\umR}{\underline{\mR}}
\newcommand{\umZ}{\underline{\mZ}}

\newcommand{\uN}{\underline{$N$}}
\newcommand{\upi}{\underline{\pi}}

\newcommand{\bxi}{\bar{\xi}}
\newcommand{\btau}{\bar{\tau}}

%shortcuts for BP<n> and such
\newcommand{\BPR}{BP_{\mathbb{R}}}
\newcommand{\0}{\langle 0 \rangle }
\newcommand{\one}{\langle 1 \rangle }
\newcommand{\two}{\langle 2 \rangle }
\newcommand{\n}{\langle n \rangle }

\newcommand{\vertsimeq}{\rotatebox{90}{$\,\simeq$}}
%%%%% formatting for quotient %%%%%

\newcommand{\bigslant}[2]{{\raisebox{.2em}{$#1$}\left/\raisebox{-.2em}{$#2$}\right.}} 

\definecolor{green'}{HTML}{009E73}
\definecolor{pink'}{HTML}{CC79A7}
\definecolor{orange'}{HTML}{D55E00}
\definecolor{darkblue'}{HTML}{0072B2}
\definecolor{yellow'}{HTML}{F0E442}
\definecolor{lightblue'}{HTML}{56B4E9}
\definecolor{gold'}{HTML}{E69F00}


% todo notes:
\usepackage[textwidth=25mm, textsize=tiny]{todonotes}

% Guchuan's notes
\newcommand{\glmar}[1]{\todo[color=green'!50!white,linecolor=green'!50!white,size=\tiny]{GL: #1}}
\newcommand{\glil}[1]{\todo[inline,color=green'!50!white,linecolor=green'!50!white,size=\normalsize]{GL: #1}}

%Jackson's Notes
\newcommand{\jmmar}[1]{\todo[color=orange'!50!white,linecolor=orange'!50!white,size=\tiny]{JM: #1}}
\newcommand{\jmil}[1]{\todo[inline,color=orange'!50!white,linecolor=orange'!50!white,size=\normalsize]{JM: #1}}

%Sarah's notes
\newcommand{\spmar}[1]{\todo[color=darkblue'!50!white,linecolor=darkblue'!50!white,size=\tiny]{SP: #1}}
\newcommand{\spil}[1]{\todo[inline,color=darkblue'!50!white,linecolor=darkblue'!50!white,size=\normalsize]{SP: #1}}

%Liz's notes
\newcommand{\ltmar}[1]{\todo[color=yellow'!50!white,linecolor=yellow'!50!white,size=\tiny]{LT: #1}}
\newcommand{\ltil}[1]{\todo[inline,color=yellow'!50!white,linecolor=yellow'!50!white,size=\normalsize]{LT: #1}}


\usepackage[style=alphabetic]{biblatex}
\usepackage{float}
\addbibresource{bib.bib}

